\documentclass{article}
\usepackage[OT1]{fontenc} % Updated from OT1 for better character rendering
\usepackage[margin=1in]{geometry}

\usepackage[style=alphabetic, backend=biber, useprefix=true]{biblatex}
% Ensure refs.bib exists in your project folder!
\addbibresource{refs.bib}
\renewcommand{\multicitedelim}{\addcomma\space}

\usepackage[dvipsnames]{xcolor}

\usepackage{hyperref}
\hypersetup{
    colorlinks = true,
    linktoc=page,
    urlcolor = BurntOrange,
    citecolor = Bittersweet,
    linkcolor = BurntOrange
}

\usepackage{enumerate}
\usepackage{subcaption}
\usepackage{amsfonts}
\usepackage{amsmath}
\usepackage{amssymb}
\usepackage{amsthm} 
\usepackage{mathtools}
\usepackage{tikz}
\usepackage{tikz-cd}
\usepackage{quiver} 

\usepackage[mathscr]{euscript}
\usepackage{scalerel}
\usepackage{wasysym}

\usepackage{titlesec}

\titleformat{\section}
  {\normalfont\Large\bfseries}{\S\thesection}{1em}{}

\setlength{\columnsep}{1cm}
\parindent = 10pt

% Definitions
\def\~#1{\mathscr{#1}}
\def\cal#1{\mathcal{#1}}
\def\To#1{\xrightarrow{#1}}

\numberwithin{equation}{section}
\numberwithin{figure}{section}

\newcommand{\figref}[1]{Figure \ref{#1}}
\newcommand{\vp}{\varphi}
\newcommand{\ve}{\varepsilon}

\newcommand{\C}{\mathbb{C}}
\newcommand{\N}{\mathbb{N}}
\newcommand{\Q}{\mathbb{Q}}
\newcommand{\R}{\mathbb{R}}
\newcommand{\Z}{\mathbb{Z}}
\newcommand{\A}{\mathbb{A}}
\newcommand{\G}{\mathbb{G}}
\newcommand{\Hom}{\mathrm{Hom}}
\newcommand{\Aut}{\mathrm{Aut}}
\newcommand{\Ext}{\mathrm{Ext}}
\newcommand{\Der}{\mathrm{Der}}
\newcommand{\GL}{\mathrm{GL}}
\newcommand{\PGL}{\mathrm{PGL}}
\newcommand{\M}{\mathrm{M}}
\newcommand{\Br}{\mathrm{Br}}
\newcommand{\Az}{\mathrm{Az}}
\newcommand{\isom}{\xrightarrow{\sim}}
\newcommand{\epi}{\twoheadrightarrow}
\newcommand{\mono}{\hookrightarrow}
\newcommand{\Spec}{\mathop{\mathrm{Spec}}}
\newcommand{\sHom}{\~H\mathrm{om}}


% Category theory symbols
\newcommand{\op}{\mathrm{op}}
\newcommand{\adj}{\dashv}
\newcommand{\dom}{\mathrm{dom}}
\newcommand{\colim}{\mathop{\mathrm{colim}}}
\newcommand{\ob}{\mathop{\mathrm{ob}}}
\newcommand{\cod}{\mathrm{cod}}
\newcommand{\coker}{\mathop{\mathrm{coker}}}
\newcommand{\im}{\mathrm{im}}
\renewcommand{\d}{\partial}
\newcommand{\id}{\mathrm{id}}
\newcommand{\res}{\mathrm{res}}
\newcommand{\et}{\acute{\mathrm{e}}\mathrm{t}}
\renewcommand{\frak}[1]{\mathfrak{#1}}
\newcommand{\sep}{\mathrm{sep}}
\newcommand{\Gal}{\mathrm{Gal}}
\newcommand{\disc}{\mathrm{disc}}


% Categories
\newcommand{\cAlg}{\mathrm{c}\~{A}\mathrm{lg}}
\newcommand{\cAlgop}{\cAlg^\op}
\newcommand{\Set}{\mathrm{Set}}
\newcommand{\Rel}{\~R\mathrm{el}}
\newcommand{\Grp}{\mathrm{Grp}}
\newcommand{\Ab}{\mathrm{Ab}}
\newcommand{\Ring}{\mathrm{Ring}}
\newcommand{\cRing}{\mathrm{c}\~{R}\mathrm{ing}}
\newcommand{\Mod}{\mathrm{Mod}}
\newcommand{\Vect}{\~V\mathrm{ect}}
\newcommand{\fdVect}{\mathrm{fd}\~{V}\mathrm{ect}}
\newcommand{\Top}{\~T\mathrm{op}}
\newcommand{\Cat}{\~C\mathrm{at}} % Fixed typo from \~Aat
\newcommand{\Mon}{\~M\mathrm{on}}
\newcommand{\Open}{\~{O}\mathrm{pen}}
\newcommand{\Rep}{\~R\mathrm{ep}}
\newcommand{\Field}{\~F\mathrm{ield}}
\newcommand{\PSh}{\mathrm{PSh}}
\newcommand{\Bilin}{\~{B}\mathrm{ilin}}
\newcommand{\Ch}{\mathrm{Ch}}
\newcommand{\Sh}{\mathrm{Sh}}
\newcommand{\Fun}{\mathrm{Fun}}
\newcommand{\Sch}{\mathrm{Sch}}
\newcommand{\Cov}{\mathrm{Cov}}
\newcommand{\LRS}{\mathrm{LRS}}


\DeclareFontFamily{U}{dmjhira}{}
\DeclareFontShape{U}{dmjhira}{m}{n}{ <-> dmjhira }{}
\DeclareRobustCommand{\yo}{\text{\usefont{U}{dmjhira}{m}{n}\symbol{"48}}}

% Theorem environments
\theoremstyle{plain}
\newtheorem{thm}{Theorem}[section]
\newtheorem*{thm*}{Theorem}
\newcommand{\thmautorefname}{Theorem}
\newtheorem{prop}[thm]{Proposition}
\newtheorem{lem}[thm]{Lemma}
\newcommand{\lemautorefname}{Lemma}
\newtheorem*{lem*}{Lemma}
\newtheorem{cor}[thm]{Corollary}
\newtheorem*{cor*}{Corollary}
\newtheorem{conj}[thm]{Conjecture}

\theoremstyle{definition}
\newtheorem{defn}[thm]{Definition}
\newtheorem*{defn*}{Definition}
\newtheorem{ex}[thm]{Example}
\AtEndEnvironment{ex}{\null\hfill$\Diamond$}
\theoremstyle{remark}
\newtheorem*{rmk}{Remark}
\newtheorem*{rmks}{Remarks}
\newtheorem*{notation}{Notation}

\renewcommand{\sectionautorefname}{\S}

\title{\'Etale cohomology}
\author{Elias Vandenberg}

\begin{document}
\maketitle
  \tableofcontents

\subsection*{Conventions and Notation}
All rings and algebras will be assumed to be commutative with identity $1 \ne 0$. The symbol $:=$ indicates that the equality in question is an definition. If $\~C$ is a category and $X$ is an object of $\~C$, we write $X \in \~C$ in place of $X \in \ob(\~C)$. We write $\Hom_{\~C}(X,Y)$ for the collection of morphisms $X \to Y$ in $\~C$. Isomorphisms between objects in a category are denoted by the symbol $\cong$, while $\simeq$ is reserved for equivalences of categories. $\Fun(\~C, \~D)$ denotes the category of functors $\~C \to \~D$. For a scheme $X$, $\Sch_{\!/\!X}$ denotes the category of schemes over $X$, i.e., the slice category over $X$, and we write $\Hom_X(S,T)$ for the hom-sets in this category.
\section{Derived functors}

We first review some basic facts about derived functors. Since we are interested in cohomology, we restrict our attention to right derived functors. The main reference for this section is Chapter 2 of Weibel's book \cite{weibel}. 

\subsection{Injective resolutions}
Throughout this section, we work in an abelian category $\~A$. Recall that a \emph{chain complex} $(A_\bullet, \d_\bullet)$ in $\~A$ is a sequence of objects and arrows in $\~A$
\[\cdots \To{\d_3} A_2 \To{\d_2} A_1 \To{\d_1} A_0 \To{\d_0} A_{-1} \To{\d_{-1}} \cdots\]
with the property that $\d_n \circ \d_{n+1} = 0$ for all $n \in \Z$. Reversing the arrows, we obtain the definition of a \emph{cochain complex} $(A^\bullet, \d^\bullet)$. Denote by $\Ch_\bullet(\~A)$ (resp. $\Ch^\bullet(\~A)$) the category of chain (resp. cochain) complexes in $\~A$. A \emph{chain homotopy} $\psi: f \simeq g$ (or $\psi: f \To{\simeq} g$) between chain maps $f,g: A_\bullet \to B_\bullet$ in $\Ch(\~A)$ is a sequence of morphisms 
\[\{(\psi_n: A_n \to B_{n + 1}) \in \~A : n \in \Z\}\]
with the property that 
\[f_n - g_n = \d^B \circ \psi_n + \psi_{n-1} \circ \d^C\]
for each $n \in \Z$.
\begin{rmk}
    In the literature, chain homotopies are often illustrated using diagrams of the form 
\[\begin{tikzcd}[column sep = scriptsize]
    \vdots && \vdots \\
    {A_{n+1}} && {B_{n+1}} \\
    {A_n} && {B_n} \\
    {A_{n-1}} && {B_{n-1}} \\
    \vdots && \vdots
    \arrow[from=1-1, to=2-1]
    \arrow[from=1-3, to=2-3]
    \arrow["{f_{n+ 1} - g_{n+1}}", from=2-1, to=2-3]
    \arrow["{\d^A}"', from=2-1, to=3-1]
    \arrow["{\d^B}", from=2-3, to=3-3]
    \arrow["{\psi_n}"{description}, from=3-1, to=2-3]
    \arrow["{f_n - g_n}", from=3-1, to=3-3]
    \arrow["{\d^A}"', from=3-1, to=4-1]
    \arrow["{\d^B}", from=3-3, to=4-3]
    \arrow["{\psi_{n-1}}"{description}, from=4-1, to=3-3]
    \arrow["{f_{n-1} - g_{n-1}}"', from=4-1, to=4-3]
    \arrow[from=4-1, to=5-1]
    \arrow[from=4-3, to=5-3]
\end{tikzcd}\]
note however that this is \emph{not} a commutative diagram. To avoid any confusion, we will refrain from using such diagrams.
\end{rmk}
\begin{defn}
    $I \in \~A$ is \emph{injective} if for each injection $f: A \mono B$ and each map $\alpha: A \to I$, there is an arrow $\beta: B \to I$ (not necessarily unique) rendering the following diagram commutative:
    \[\begin{tikzcd}
    0 & A & B \\
    && I
    \arrow[from=1-1, to=1-2]
    \arrow["f", hook, from=1-2, to=1-3]
    \arrow["\alpha"', from=1-2, to=2-3]
    \arrow["{\exists\beta}", from=1-3, to=2-3]
\end{tikzcd}\]
Equivalently, $I$ is injective if the contravariant hom-functor $X \mapsto \Hom_{\~A}(X, I)$ is exact.
\end{defn}

\begin{defn}
    $\~A$ is said to have \emph{enough injectives} if for every object $A \in \~A$ there is an injection $A \mono I$ with $I \in \~A$ injective.  
\end{defn}

\begin{defn}
    A \emph{right resolution} of $A \in \~A$ is a cochain complex $I^\bullet$ with $I^i = 0$ for $i < 0$ and a map $A \to I^0$ such that the complex 
    \[0 \to A \to I^0 \to I^1 \to I^2 \to \cdots\]
    is exact. An \emph{injective resolution} of $A \in \~A$ is a right resolution of $A$ with the property that each $I^i$ is injective.
\end{defn}
A useful result will be the following:
\begin{lem}
    If $\~A$ has enough injectives, then every object in $\~A$ has an injective resolution.
\end{lem}

\subsection{Definition of derived functors}
Let $\~A$ be an abelian category with enough injectives, and suppose $F: \~A \to \~B$ is a left exact functor into another abelian category $\~B$. Recall this means that for each exact sequence $0 \to A \to B \to C \to 0$ in $\~A$, the sequence
\[0 \to F(A) \to F(B) \to F(C)\]
is exact in $\~B$. The main issue with such a sequence is, of course, the failure of exactness on the right. Right derived functors arise as a means of remedying this issue. Given our left-exact functor $F: \~A \to \~B$, our goal is to ``derive'' a series of functors $R^iF : \~A \to \~B$ with the following two properties:
\begin{enumerate}
    \item $R^0F = F$
    \item Given a short exact sequence $0 \to A \to B \to C \to 0$ in $\~A$, applying the functors $R^iF$ yields a long exact sequence
    \begin{equation}\label{eq:right_derived_les1}\begin{tikzcd}[sep = scriptsize]
    0 & F(A) & F(B) & F(C) \\
    & {R^1F(A)} & {R^1F(B)} & {R^1F(C)} \\
    & {R^2F(A)} & \cdots
    \arrow[from=1-1, to=1-2]
    \arrow[from=1-2, to=1-3]
    \arrow[from=1-3, to=1-4]
    \arrow[from=1-4, to=2-2]
    \arrow[from=2-2, to=2-3]
    \arrow[from=2-3, to=2-4]
    \arrow[from=2-4, to=3-2]
    \arrow[from=3-2, to=3-3]
\end{tikzcd}\end{equation}
\end{enumerate}

\begin{defn} 
Suppose $F: \~A \to \~B$ is an exact functor between abelian categories with $\~A$ having enough injectives. Given $A \in \~A$, take an injective resolution
\[I_A^\bullet : ~ 0 \to A \to I^0 \To{\d} I^1 \To{\d} I^2 \To{\d} \cdots\]
and apply $F$ to get a new complex
\[F(I_A^\bullet) : ~ 0 \to F(A) \to F(I^0) \To{F(\d)} F(I^1) \To{F(\d)} F(I^2) \To{F(\d)} \cdots\]
We then define the \emph{right derived functors} $R^i F$ ($i \ge 0$) of $F$ by 
\begin{equation}R^iF(A) := H^i(F(I_A^\bullet)) = \ker\left(F\left(\d^i\right)\right)/~\im\left(F(\d^{i-1})\right).\end{equation}
\end{defn}
From this definition, it is not immediately clear that the functors $R^iF$ are well-defined: the definition depends on a \emph{choice} of injective resolution $A \to I^\bullet$, and it is not at all obvious that two different injective resolutions of $A$ will give rise to the same right-derived functor. To prove that the functors $R^iF$ are well-defined, we will need two additional results:
\begin{lem}\label{lem:chain_homotopy_lem}
    Chain homotopic maps induce the same maps on (co)homology.
\end{lem}
% \begin{proof}
%     It suffices to prove the result for homology. Suppose $(A_\bullet,  B_\bullet \in \Ch_\bullet(\~A)$, and take 
% \end{proof}
\begin{thm}[Comparison theorem, \cite{weibel} Theorem 2.3.7]\label{thm:comparison}
     Let $\~A$ be an abelian category with enough injectives, and take $B \in \~A$ and an injective resolution $B \to I^\bullet$. Suppose we have a morphism $f': A \to B$ in $\~A$. Then for every resolution $A \to J^\bullet$, there is a cochain map $f: J^\bullet \to I^\bullet$ lifting $f'$, as depicted in the following diagram:
     \[\begin{tikzcd}
    0 & A & {J^0} & {J^1} & {J^2} & \cdots \\
    0 & B & {I^0} & {I^1} & {I^2} & \cdots
    \arrow[from=1-1, to=1-2]
    \arrow[from=1-2, to=1-3]
    \arrow["{f'}"', from=1-2, to=2-2]
    \arrow[from=1-3, to=1-4]
    \arrow["\exists"', from=1-3, to=2-3]
    \arrow[from=1-4, to=1-5]
    \arrow["\exists"', from=1-4, to=2-4]
    \arrow[from=1-5, to=1-6]
    \arrow["\exists"', from=1-5, to=2-5]
    \arrow[from=2-1, to=2-2]
    \arrow[from=2-2, to=2-3]
    \arrow[from=2-3, to=2-4]
    \arrow[from=2-4, to=2-5]
    \arrow[from=2-5, to=2-6]
\end{tikzcd}\]
Moreover, $f$ is unique up to homotopy equivalence.
\end{thm}

We may now easily deduce that the right-derived functors of $F$ are well-defined.
\begin{cor}
    Suppose $F: \~A \to \~B$ is an exact functor between abelian categories with $\~A$ having enough injectives. Then the functors $R^iF$ are well-defined.
\end{cor}
\begin{proof}
Take two injective resolutions $I = (A \to I^\bullet)$ and $J = (A \to J^\bullet)$ of $A \in \~A$. By Theorem \ref{thm:comparison} above, we can find a chain map $f: I^\bullet \to J^\bullet$, unique up to homotopy equivalence, which lifts $\id_A: A \to A$:
     \[\begin{tikzcd}
    0 & A & {I^0} & {I^1} & {I^2} & \cdots \\
    0 & A & {J^0} & {J^1} & {J^2} & \cdots
    \arrow[from=1-1, to=1-2]
    \arrow[from=1-2, to=1-3]
    \arrow["{\id_A}"', from=1-2, to=2-2]
    \arrow[from=1-3, to=1-4]
    \arrow["\exists"', from=1-3, to=2-3]
    \arrow[from=1-4, to=1-5]
    \arrow["\exists"', from=1-4, to=2-4]
    \arrow[from=1-5, to=1-6]
    \arrow["\exists"', from=1-5, to=2-5]
    \arrow[from=2-1, to=2-2]
    \arrow[from=2-2, to=2-3]
    \arrow[from=2-3, to=2-4]
    \arrow[from=2-4, to=2-5]
    \arrow[from=2-5, to=2-6]
\end{tikzcd}\]
As $f, g: I^\bullet \to J^\bullet$ are chain homotopic, so are $F(f), F(g): F(I^\bullet) \to F(J^\bullet)$, therefore by Lemma \ref{lem:chain_homotopy_lem} we have isomorphisms 
\[H^i(F(I)) \cong H^i(F(J)),\]
so $R^iF(A)$ is independent of the injective resolution chosen.
\end{proof}

We now fix an abelian category $\~A$ with enough injectives. To prove that the right derived functors induce the desired long exact sequence (\ref{eq:right_derived_les1}), we will need the following three results from homological algebra:
\begin{lem}\label{lem:inj_proj_ses}
    Let $0 \to A \to B \to C \to 0$ be a short exact sequence in $\~A$. If $A$ is injective, then the sequence splits. Dually, if $C$ is projective, then the sequence splits. 
\end{lem}
\begin{lem}\label{lem:les_on_cohomology}
    Given a short exact sequence of complexes in $\Ch^\bullet(\~A)$
    \[0 \to A^\bullet \to B^\bullet \to C^\bullet \to 0\]
    there is a long exact sequence in cohomology
    \[\cdots \to H^{n-1}(B^\bullet) \to H^{n-1}(C^\bullet) \to H^n(A^\bullet) \to H^n(B^\bullet) \to \cdots\]
\end{lem}

\begin{lem}[Horseshoe lemma]\label{lem:horseshoe}
    Suppose $\~A$ has enough injectives, and take a short exact sequence $0 \to A \to B \to C \to 0$ in $\~A$. Let $A \to I^\bullet$, $C \to J^\bullet$ be two injective resolutions. Then $B \to I^\bullet \oplus J^\bullet$ defines an injective resolution making the following diagram commute:
    \[\begin{tikzcd}
    0 & A & B & C & 0 \\
    0 & {I^\bullet} & {I^\bullet \oplus J^\bullet} & {J^\bullet} & 0
    \arrow[from=1-1, to=1-2]
    \arrow[from=1-2, to=1-3]
    \arrow[from=1-2, to=2-2]
    \arrow[from=1-3, to=1-4]
    \arrow[from=1-3, to=2-3]
    \arrow[from=1-4, to=1-5]
    \arrow[from=1-4, to=2-4]
    \arrow[from=2-1, to=2-2]
    \arrow[from=2-2, to=2-3]
    \arrow[from=2-3, to=2-4]
    \arrow[from=2-4, to=2-5]
\end{tikzcd}\]
\end{lem}
\begin{thm}\label{eq:right_derived_les}
    Given a short exact sequence
    \[0 \to A \to B \to C \to 0\]
    in $\~A$, there is a long exact sequence
    \begin{equation} \cdots \to R^{i-1}F(C) \to R^iF(A) \to R^iF(B) \to R^iF(C) \to R^{i + 1}F(A) \to \cdots \end{equation}
\end{thm}
\begin{proof}
    Let $0 \to A \To{f} B \To{g} C \to 0$ be an exact sequence and take injective resolutions $A \to I_A^\bullet$ and $C \to I_C^\bullet$. Using Lemma \ref{lem:horseshoe} above, we can produce an injective resolution $B \to I_B^\bullet$ fitting into the commuting diagram
        \[\begin{tikzcd}
    0 & A & B & C & 0 \\
    0 & {I_A^\bullet} & {I_B^\bullet} & {I_C^\bullet} & 0
    \arrow[from=1-1, to=1-2]
    \arrow[from=1-2, to=1-3]
    \arrow[from=1-2, to=2-2]
    \arrow[from=1-3, to=1-4]
    \arrow[from=1-3, to=2-3]
    \arrow[from=1-4, to=1-5]
    \arrow[from=1-4, to=2-4]
    \arrow[from=2-1, to=2-2]
    \arrow[from=2-2, to=2-3]
    \arrow[from=2-3, to=2-4]
    \arrow[from=2-4, to=2-5]
\end{tikzcd}\]
For an integer $n$, consider the exact sequence $0 \to I_A^n \to I_B^n \to I_C^n \to 0$. As $I_A^n$ is injective, Lemma \ref{lem:inj_proj_ses} tells us that the sequence splits for every $n$: $I^n_B \cong I_A^n \oplus I_C^n$. Moreover, $F$ is exact, so in particular it is additive, therefore the sequence $0 \to F(I_A^n) \to F(I_B^n) \to F(I_C^n) \to 0$ is also split exact for every $n$. This gives a short exact sequence of complexes
\[0 \to F(I_A^\bullet) \to F(I_B^\bullet) \to F(I_C^\bullet) \to 0\]
which by Lemma \ref{lem:les_on_cohomology} induces a long exact sequence
\[\cdots \to H^{n-1}(F(I_B^\bullet)) \to H^{n-1}(F(I_C^\bullet)) \to H^n(F(I_A^\bullet)) \to H^n(F(I_B^\bullet)) \to \cdots\]
Since $R^nF(A) = H^n(F(I_A^\bullet))$, we obtain the long exact sequence (\ref{eq:right_derived_les}).
\end{proof}
Our next definition describes the prototypical example of right derived functors  
\begin{defn}
    For each $A \in \Mod_R$, the functor $\Hom_R(A,-)$ is left exact. Its right derived functors are called the \emph{Ext-groups}:
    \[\Ext^i_R(A,B) := R^i \Hom_R(A, -)(B).\]
\end{defn}

\begin{ex}
Let $X$ be a topological space and denote by $\Sh_{\Ab}(X)$ the category of sheaves of abelian groups on $X$. Recall that the \emph{global sections functor} $\Gamma(X, -): \Sh_{\Ab}(X) \to \Ab$ is the functor defined by $\Gamma(X, \~F) := \~F(X)$ for each $\~F \in \Sh_\Ab(X)$. Taking global sections is right adjoint to the constant sheaf functor $A \mapsto \underline{A}$:
\[\Hom_{\Sh(X)}(\underline{A}, \~F) \cong \Hom_{\Ab}(A, \Gamma(X,\~F)).\]
In particular, $\Gamma(X, -)$ is left-exact: given a short exact sequence $0 \to \~F \to \~G \to \~H \to 0$ on the same space $X$, the sequence
\[0 \to\Gamma(X, \~F) \to\Gamma(X, \~G) \to \Gamma(X, \~H) \]
is exact. To address the failure of right-exactness, we will compute the right derived functors $R^i\Gamma(X, \~F)$, which are called the \emph{cohomology functors of the sheaf $\~F$}. We write $H^i(X, \~F) := R^i\Gamma(X, \~F)$.
\end{ex}


\section{The \'etale site and cohomology}

% \subsection{Differentials}

% For the sake of completeness, we use this section to briefly review the definition of K\"ahler differentials. Let $f: R \to A$ be a ring homomorphism and $M$ an $A$-module. A map of $R$-modules $D: A \to M$ is called an \emph{$R$-linear derivation} (or simply a \emph{derivation}) if it annihilates elements in $f(R)$ and satisfies the \emph{Leibniz rule}:
% \begin{equation}
%     D(ab) = aD(b) + bD(a) \quad \text{for all } a, b \in A
% \end{equation}
% The set $\Der_R(A, M)$ of all $R$-linear derivations $A \to M$ is itself an $A$-module: For $D, D' \in \Der_R(A, M)$, define $D + D': A \to M$ by $a \mapsto D(a) + D'(a)$, and $cD: A \to M$ by $a \mapsto cD(a)$ for all $c \in A$.
% Given $R \to A$ as above, the \emph{module of K\"ahler differentials} $\Omega_{A/R}$ is defined as the free $A$-module of the symbols $\{d(a) : a \in A \}$, subject to the following relations for all $a, b \in A$ and all $r \in R$:
% \begin{enumerate}
%     \item $d(a + b) = da + db$
%     \item $d(0) = 0$
%     \item $d(ra) = r (da)$ 
%     \item $d(ab) = a(db) + b(da)$ (\emph{Leibniz rule})
%     \item $d(1) = 0$
% \end{enumerate}
% The map $d: A \to \Omega_{A/R}$ given by $d: a \mapsto d(a)$ is a derivation, called the \emph{universal $R$-linear derivation}. Observe that this map is indeed $R$-linear. We often write $da$ in place of $d(a)$.

% \begin{rmk}
%     While the definition of $\Omega_{A/R}$ works for any ring homomorphism $R \to A$, it is often helpful to restrict our attention to the case where $A$ is itself an $R$-algebra (so the image of $R \to A$ lies in the centre of $A$).
% \end{rmk}

% \begin{ex}
% Consider $R[x]$. In theory, there is a generator $dp(x)$ for every $p(x) \in R[x]$, but using our relations, we can reduce things. For instance, $x^2$ can be reduced as follows
% \[d(x^2) = d(x \cdot x) = xdx + x dx = 2xdx\]
% and in general, we have $d(x^n) = nx^{n - 1}dx$. Using the module structure, we can then deduce that any $dp(x)$ reduces to a linear combination of $dx$, so the $\Omega_{R[x]/R}$ is the free $R[x]$-module generated by the single symbol $dx$. More generally, we have that $\Omega_{R[x_1, \cdots, x_n]/R}$ is the free $R[x_1, \cdots, x_n]$-module generated by $(dx_1, \cdots, dx_n)$:
% \[\Omega_{R[x_1, \cdots, x_n]/R} = \bigoplus^n_{i = 1}R[x_1, \cdots, x_n] dx_i\]
% \end{ex}

\subsection{Review of Grothendieck topologies}
In this section, we briefly review the theory of Grothendieck sites as described in \cite{SGL}. Suppose $\~C$ is a small category, and denote by $\PSh(\~C) = \Fun(\~C^\op, \Set)$ the category of presheaves on $\~C$. We write $\yo: \~C \to \PSh(\~C)$ for the Yoneda embedding ($\yo(A) = \Hom_{\~C}(-, A)$ for all $A \in \~C$). 
\begin{defn}
    A \emph{sieve} $S$ on an object $A \in \~C$ is a subfunctor $S \subseteq \yo(A)$. More concretely, a sieve is a family $S$ of morphisms in $\~C$, all of which have codomain $A$, such that for all $f: B \to A$ in $S$ and all arrows $g: C \to B$, we have $f \circ g \in S$. 
\end{defn}
Hence a sieve is a collection of arrows with codomain $A$ which is closed under precomposition. Sieves therefore play a similar role to right ideals in ring theory. If $S$ is a sieve on $A \in \~C$ and $h: B \to A$ is any arrow in $\~C$, we can define a sieve $h^\ast(S)$ on $B$ by 
\[h^\ast(S) := \left\{g : \cod(g) = B, \, h\circ g \in S \right\}.\]
Moreover, for any object $A \in \~C$, the collection of all possible maps into $A$ clearly defines a sieve $t_A$, called the \emph{maximal sieve}:
\[t_A := \left\{f : \cod(f) = A\right\}.\]
Notice that if $S$ is a sieve on $A \in \~C$, we have 
\[ S = t_A \iff 1_A \in S, \]
this is another instance of the similarity between sieves and right ideals.

\begin{defn}\label{defn:GT}
    A \emph{Grothendieck topology} on a category $\~C$ is a function $J$ which assigns to each object $A \in \~C$ a collection $J(A)$ of sieves on $\~C$, such that:
    \begin{enumerate}
        \item \label{GT:1} The maximal sieve $t_A$ is in $J(A)$. 
        \item \label{GT:2}(\emph{Stability axiom}) If $S \in J(A)$, then $h^\ast(S) \in J(B)$ for any arrow $h: B \to A$.   
        \item \label{GT:3}(\emph{Transitivity axiom}) If $S \in J(A)$ and $R$ is a sieve on $A$ such that $f^*(R) \in J(B)$ for all $f: B \to A$ in $S$, then $R \in J(A)$.
    \end{enumerate}
    A \emph{site} is a pair $(\~C, J)$ consisting of a small category $\~C$ and a Grothendieck topology $J$ on $\~C$.
\end{defn}

Suppose $(\~C, J)$ is a site. Then if $S \in J(A)$ for some $A \in \~C$, any larger sieve $R \supseteq S$ on $A$ must also lie in $J(A)$. To see this, suppose we have an arrow $h: B \to A$ in $S$. Then certainly $1_B \in h^*(S)$, so $h^*(S) = t_B$, the maximal sieve on $B$. As $h^*(S) \subseteq h^*(R)$, we must also have $h^*(R) = t_B$, hence $h^*(R) \in J(D)$ by Definition \ref{defn:GT}. This is true for all $h: B \to A$ in $S$, so transitivity gives $R \in J(A)$.

We can rephrase the axioms of a Grothendieck topology using the more suggestive terminology of ``coverings''. If $S \in J(A)$, we call $S$ a \emph{covering sieve} and say that $S$ \emph{covers} $A$. We say that $S$ \emph{covers an arrow $f: B \to A$} if $f^*(S)$ covers $B$ (therefore $S$ covers $A$ if and only if $S$ covers the identity arrow on $A$). The axioms for a Grothendieck topology may then be phrased as follows:
\begin{itemize}
    \item[1*.] If $S$ is a sieve on $A$ and $f \in S$, then $S$ covers $f$. 
    \item[2*.] If $S$ covers $F: B \to A$, then $S$ also covers $f \circ g$ for all $g: C \to B$.
    \item[3*.] If $S$ covers $f: B \to A$ and $R$ is a sieve on $A$ which covers all arrows of $S$, then $R$ covers $f$. 
\end{itemize}
More formally we make the following definition:
\begin{defn}\label{defn:GT2}
    A \emph{Grothendieck} topology on a category $\~C$ consists of a collection $\Cov(\~C)$ of families of morphisms $\{U_i \to U\}_{i \in I}$ called \emph{coverings of $\~C$}, satisfying the following properties:
    \begin{enumerate}
        \item If $V \to U$ is an isomorphism, then $\{V \to U\} \in \Cov(\~C)$.
        \item If $\{U_i \to U\}_{i \in I} \in \Cov(\~C)$ and for each $i$ we have $\{V_{ij} \to U_i\}_{j \in J_i} \in \Cov(\~C)$, then $\{V_{ij} \to U\}_{i \in I, j \in J_i} \in \Cov(\~C)$.
        \item If $\{U_i \to U\}_{i \in I} \in \Cov(\~C)$ and $V \to U$ is an arrow in $\~C$, then $U_i \times_{U} V$ exists for all $i$ and $\{U_i \times_{U} V \to V\}_{i \in I} \in \Cov(\~C)$.
    \end{enumerate}
\end{defn}
This definition of a Grothendieck topology is the one most frequently encountered in algebraic geometry.

\subsection{\'Etale morphisms}
The standard Zariski topology of a scheme is often far too coarse. To remedy this issue, we would like to equip our schemes with a finer topology. In the case of a variety over $\C$, we identify \'etale morphisms with local diffeomorphisms in the topological sense (for instance, covering space maps). These morphisms give rise to a topology (the \'etale topology) which is much finer than the Zariski topology. In this topology, cohomological invariants become much more well-behaved. The main reference for this section is \cite{stacks}.

Throughout this section, we will make frequent reference to \emph{flat} morphisms of schemes. A module $M$ over a ring $R$ is called \emph{flat} if the tensor product functor $(-)\otimes_R M$ is exact, and a ring morphism $R \to A$ is \emph{flat} if $A$ is a flat $R$-module under the action induced by this homomorphism. A morphism $f: X \to Y$ of ringed spaces is \emph{flat at $x \in X$} if the associated ring homomorphism $\~O_{Y, f(x)} \to \~O_{X, x}$ is flat, and we say that $f$ is \emph{flat} if it is flat at every $x \in X$. When $\~F$ is an $\~O_X$-module, we say that $\~F$ is \emph{flat} if the functor 
\[
   (-)\otimes_{\~O_X}\~F: \Mod_{\~O_X} \to \Mod_{\~O_X}, \quad \~G \mapsto \~G \otimes_{\~O_X} \~F
\]
is exact (for more details, see \S3.9 in \cite{hartshorne}). 

Recall that a morphism of rings $R \to A$ is said to be of \emph{finite type} (or that $A$ is a \emph{finite type $R$-algebra}) if there is $n \in \Z_{\ge 1}$ and a surjection of $R$-algebras $R[x_1, \cdots, x_n] \epi A$. $R \to A$ is of \emph{finite presentation} if there are integers $n, m \in \Z_{\ge 1}$ and polynomials $f_j$ such that $A \cong R[x_1, \hdots, x_n]/\langle f_1, \hdots, f_m \rangle$.

\begin{notation}
    For a ringed space $(X, \~O_X)$, we denote by $\frak{m}_x$ the maximal ideal of the local ring $\~O_{X,x}$ and $\kappa(x) := \~O_{X,x}/\frak{m}_x$ the residue field of $X$ at $x$.
\end{notation}

\begin{defn}[\cite{stacks}, Definition \href{https://stacks.math.columbia.edu/tag/01TP}{01TP}]
    A morphism of schemes $f: X \to Y$ is said to be of \emph{finite presentation} at $x \in X$ if there is an affine open neighbourhood $\Spec(B) = U \subseteq X$ of $x$ and an affine open $\Spec(A) = V \subseteq Y$ with $f(U) \subseteq V$ such that the induced ring map $A \to B$ is of finite presentation. We say $f$ is \emph{locally of finite presentation} if it is of finite presentation at every $x \in X$.
\end{defn}

\begin{defn}
    A morphism of rings $A \to B$ is \emph{unramified} if it is of finite type and $\Omega_{B/A} = 0$. 
\end{defn}
\begin{defn}[\cite{stacks}, Definition \href{https://stacks.math.columbia.edu/tag/02G4}{02G4}]
 A morphism of schemes $f: X \to Y$ is said to be \emph{unramified at $x \in X$} if there exists an affine open neighbourhood $\Spec(B) = U \subseteq X$ of $x$ and an affine open $\Spec(A) = V \subseteq Y$ with $f(U) \subseteq V$ such that the induced ring map $A \to B$ is unramified. $f$ is \emph{unramified} if it is unramified at every point in $X$.
\end{defn}

\begin{defn}[\cite{stacks}, Definition \href{https://stacks.math.columbia.edu/tag/00U1}{00U1}]
    A morphism of rings $A \to B$ is called \emph{\'etale} if it is smooth (in the sense of \cite{stacks}, Definition \href{https://stacks.math.columbia.edu/tag/00T2}{00T2}) and $\Omega_{B/A} = 0$.
\end{defn}

\begin{lem}[\cite{stacks}, Lemma \href{https://stacks.math.columbia.edu/tag/00U9}{00U9}]
    A morphism of rings $A \to B$ is \'etale if and only if 
    \[B \cong A[x_1, \hdots, x_n]/\langle f_1, \hdots, f_n \rangle\]
    such that $\Delta = \det\left(\frac{\partial f_i}{\partial x_j}\right)$ is invertible in $B$.
\end{lem}

\begin{defn}\label{def:etale}
    A morphism of schemes $(f, f^\sharp): (X, \~O_X) \to (Y, \~O_Y)$ is called \emph{\'etale} if it satisfies any one of the following equivalent properties:
    \begin{enumerate}
        \item $f$ is flat and unramified. 
        \item $f$ is smooth and unramified.
        \item $f$ is smooth and $\Omega_{X/Y} = 0$.
        \item \label{et3} $f$ is flat, locally of finite presentation, and for every $y \in Y$, the fibre $f^{-1}(y)$ is a disjoint union of points, each of which is the spectrum of finite separable extension of the residue field $\kappa(y)$.
        \item For every point $x$ of $X$, there is an affine open neighbourhood $\Spec(S) = U \subseteq X$ of $x$ and an affine open $\Spec(R) = V \subseteq Y$ with $f(U) \subseteq V$ such that the induced ring map $R \to S$ is \'etale. 
        \item For all affine opens $U \subseteq X$, $V \subseteq Y$ with $f(U) \subseteq V$, the morphism of rings $\~O_Y(V) \to \~O_X(U)$ is \'etale.
        \item There is an open covering $Y = \bigcup_{j \in J} V_j$ and coverings $f^{-1}(V_j) = \bigcup_{i \in I_j}U_i$ such that $f|_{U_i}: U_i \to V_j$ is \'etale for each $j \in J$, $i \in I_j$.
        \item There is a covering by open affines $Y = \bigcup_{j \in J} V_j$ and affine open coverings $f^{-1}(V_j) = \bigcup_{i \in I_j} U_i$, such that the composite
        \[\~O_Y(V_j) \To{f_{V_j}^\sharp} \~O_X(f^{-1}(V_j)) \To{\res_{f^{-1}(V_j), U_i}} \~O_X(U_i)\]
        is \'etale for each $j \in J$, $i \in I_j$.
    \end{enumerate}
\end{defn}

\begin{rmk}
    The third characterization of \'etale maps described above is (in my opinion) the most confusing, so it worth explaining this in greater detail. Let $f: (X, \~O_X) \to (Y, \~O_Y)$ be a morphism of schemes. Take a point $y \in Y$ and an open neighbourhood $U \subseteq Y$ of $y$ with $U \cong \Spec(R)$. Then $y$ corresponds to a unique prime ideal $\frak{p}$ with $\~O_{Y,y} = R_{\frak{p}}$ and $\kappa(y) = R_{\frak{p}}/\frak{m}_y$, and the composite map 
    \[R \to R_{\frak{p}} \to \kappa(y)\]
    induces a morphism of schemes $i_y: \Spec(\kappa(y)) \to Y$. In this situation, the fibre of $f$ over $y$ is given by the pullback $f^{-1}(y) = X \times_Y \text{Spec}(\kappa(y))$:
\[\begin{tikzcd}
    {f^{-1}(y)} & X \\
    {\Spec(\kappa(y))} & Y
    \arrow[from=1-1, to=1-2]
    \arrow[from=1-1, to=2-1]
    \arrow["\lrcorner"{anchor=center, pos=0.125}, draw=none, from=1-1, to=2-2]
    \arrow["f", from=1-2, to=2-2]
    \arrow["{{i_y}}"', from=2-1, to=2-2]
\end{tikzcd}\]
Definition \ref{def:etale}(3) says that this fibre can be written as a disjoint union
\[f^{-1}(y) \cong \Spec(L_1) \sqcup \Spec(L_2) \sqcup \cdots \sqcup \Spec(L_n)\]
with each $L_i$ a finite separable field extension of $\kappa(y)$. In terms of global sections, it means there is a ring isomorphism
\[\Gamma \left(f^{-1}(y), \~O_{f^{-1}(y)} \right) \cong L_1 \times L_2 \times \cdots \times L_n\]
(this is a direct consequence of the isomorphism $\Gamma(X \sqcup Y, \~{O}) \cong \Gamma(X, \~{O}) \times \Gamma(Y, \~{O})$).
\end{rmk}

Our next lemma lists some of the basic properties of \'etale morphisms. These results can all be found in Sections \href{https://stacks.math.columbia.edu/tag/02GH}{02GH} and \href{https://stacks.math.columbia.edu/tag/03PA}{03PA} of \cite{stacks}; we have compiled them into a single lemma for the sake of convenience.
\begin{lem}\label{lem:etale_properties}
Let $X, Y$, and $S$ be schemes. 
\begin{enumerate}
     \item If $f: X \to Y$ is \'etale, then for all open subschemes $U \subseteq X$, $V \subseteq Y$ with $f(U) \subseteq V$, the restriction $f|_U: U \to V$ is \'etale.
     \item The composition of two \'etale morphisms is \'etale, and if $X \to S$ and $Y \to S$ are \'etale, then any $S$-morphism $X \to Y$ is \'etale.
     \item The base change of an \'etale morphism is \'etale.
     \item Fibred products of \'etale morphisms are \'etale.
     \item Any open immersion is \'etale.
     \item An \'etale morphism is open.
     \item An \'etale morphism is flat.
\end{enumerate}
\end{lem}
\subsection{\'Etale topology and sheaves}

\begin{defn}[\cite{stacks}, Definition \href{https://stacks.math.columbia.edu/tag/0215}{0215}]
    An \emph{\'etale covering} of a scheme $X$ is a family $\{f_i: U_i \to X\}_{i \in I}$ of morphisms of schemes such that each $f_i$ is \'etale and $X = \bigcup_{i \in I} f_i(U_i)$.
\end{defn}

\begin{lem}[\cite{stacks}, Lemma \href{https://stacks.math.columbia.edu/tag/0216}{0216}]
Any Zariski covering is an \'etale covering.
\end{lem}
\begin{proof}
    By definition, a Zariski covering of a scheme $X$ is given by a collection $\{f_i: U_i \to X\}_{i \in I}$ of open immersions, each of which is \'etale by Lemma \ref{lem:etale_properties}(4). This is enough to conclude that $\{f_i: U_i \to X\}_{i \in I}$ gives an \'etale covering of $X$.
\end{proof}

We now show that the collection of \'etale coverings constitutes a Grothendieck topology (in the sense of Definition \ref{defn:GT2}):

\begin{prop}[\cite{stacks} Lemma \href{https://stacks.math.columbia.edu/tag/0217}{0217}]
    Let $X$ be a scheme.
    \begin{enumerate}
        \item If $X' \to X$ is an isomorphism, then $\{X' \to X\}$ is an \'etale covering of $X$. 
        \item If $\{U_i \to X\}_{i \in I}$ is an \'etale covering of $X$ and for each $i$ we are given an \'etale covering $\{U_{ij} \to U_i\}_{j \in J_i}$, then $\{U_{ij} \to X\}_{i \in I, j \in J_i}$ is an etale covering. 
        \item If $\{U_i \to X\}_{i \in I}$ is an \'etale covering and $X \to Y$ is a morphism of schemes, then $\{X \times_Y U_i \to Y\}_{i \in I}$ is an \'etale covering.
    \end{enumerate}
\end{prop}

\begin{defn}
    For a scheme $X$, we define the \emph{\'etale topology} on $\Sch_{/X}$ to be the Grothendieck topology whose coverings are given by \'etale coverings. The \emph{big \'etale site} $\left(\Sch_{/X}\right)_{\et}$ on $X$ is the slice category $\Sch_{/X}$ equipped with the \'etale topology. The \emph{small \'etale site} $X_{\et} \mono \left(\Sch_{/X}\right)_{\et}$ is the full subcategory of $\Sch_{/X}$ consisting of those schemes $U$ over $X$ with $U \to X$ \'etale. By a slight abuse of notation, we often write $U \in X_{\et}$ in place of $(U \to X) \in X_{\et}$.
\end{defn}

Hence $X_{\et}$ can be viewed as the category whose objects are \'etale morphisms and whose arrows are \'etale morphisms over $X$, where $\{f_i: U_i \to U\}$ covers $U$ if $\coprod_i U_i \epi U$ is surjective.

\begin{rmk}[Size issues \frownie]
    We have been slightly cavalier with our definition of the big \'etale site. Recall that when defining a Grothendieck site, the underlying category is usually assumed to be small. However, $\Sch_{/X}$ is certainly not a small category, so we must take care when defining a Grothendieck topology in this context. The way to fix this is to choose an ordinal $\alpha$ and consider the full subcategory $\Sch_{/X}^\alpha \mono \Sch_{/X}$ of $\alpha$-bounded subschemes (as defined in \href{https://stacks.math.columbia.edu/tag/000H}{\S000H} of \cite{stacks}). Heuristically, this means we consider only those schemes $(U, \~O_U) \in \Sch_{/X}$ with $|U| \le \alpha$ and $|\~O_U(V)| \le \alpha$ for all $V \in X_{\et}$. The big \'etale site can then be formally defined by equipping $\Sch_{/X}^\alpha$ with the \'etale topology. It is also not immediately obvious why $X_{\et}$ is a site; a full proof of this can be found in \cite{stacks}, Lemma \href{https://stacks.math.columbia.edu/tag/021C}{021C} (the proof boils down to the fact that \'etale maps are locally of finite presentation). We will mostly work with $X_{\et}$ in these notes.
\end{rmk}

\begin{defn}
    A \emph{\'etale sheaf} is a sheaf on the site $X_{\et}$. In full detail, this means an \'etale sheaf is a presheaf $\~F: X_{\et}^\op \to \Set$ such that for any $U \in \Sch_{/X}$ and any cover $\{U_i \to U\}_{i \in I}$ of $U$, the diagram
    \[\~F(U) \to \prod_i \~F(U_i) \rightrightarrows \prod_{i,j} \~F(U_i \times_X U_j)\]
    is an equalizer. We write $\Sh_{\Ab}(X_{\et})$ for the category of abelian sheaves on $X_{\et}$. We have a \emph{global sections functor}:  
    \[\Gamma(X, -): \Sh_{\Ab}(X_{\et}) \to \Ab, \quad \~F \mapsto \~F(X)\]
    which is left-exact. The category $\Sh(X_{\et})$ of all sheaves on the site $X_{\et}$ is called the \emph{\'etale topos}.
\end{defn}

\begin{rmk}
$\Sh(X_{\et})$ has all products, and has a terminal object given by $U \mapsto \{\varnothing\}$ for all $U \in X_{\et}$ and an initial object given by the functor $\Hom_{X_{\et}}(-, X)$. $\Sh(X_{\et})$ and $\Sh_{\Ab}(X_{\et})$ have the same terminal object, but their initial objects are not the same. As in ordinary sheaf theory, any $\~F \in \Sh(X_{\et})$ converts disjoint unions into products: $\~F\left(\coprod_i U_i\right) = \prod_i \~F(U_i)$. The topos $\Sh_{\Ab}(X_{\et})$ is an abelian category with enough injectives, so it makes sense to discuss cohomology in this category.
\end{rmk}

\begin{ex}
    Let $\~F \in \Sh(X_{\et})$, and consider the empty cover of the empty object $\varnothing$. The product over an empty collection of sets given by $\{\varnothing\}$, so we deduce that $\~F(\varnothing) = \{\varnothing\}$, the final object in $\Set$.
\end{ex}

\begin{ex}
    Given a scheme $U$ over another scheme $X$, Yoneda's lemma tells us that the functor 
    \[X_{\et} \to \Sh(X_{\et}), \quad U \mapsto h_U := \Hom_X(-, U)\]
    is fully faithful. For any $T \to X$ in $X_{\et}$, the \emph{representable sheaf} $h_U: X_{\et}^\op \to \Set$ is defined by 
        \[h_U(T) := \Hom_{S}(T,U)\]
    we say that $\~F \in \Sh(X_{\et})$ is \emph{representable} if there is exists $U \in \Sch_{/X}$ and an isomorphism of sheaves $h_U \cong \~F$. 
\end{ex}

\begin{ex}
    We have functors $\G_a: X_{\et}^\op \to \Ab$ and $\G_m: X_{\et}^\op \to \Ab$ given by
    \[\G_a(U) := \Gamma(U, \~O_U) \quad \text{and} \quad \G_m(U) := \Gamma(U, \~O_U)^\times\]
    for each $U \in X_{\et}$. $\G_a$ and $\G_m$ define \'etale sheaves on $X_{\et}$ (see \S1.2.7 of \cite{conrad} and \S1.6 of \cite{milne} for more details). For $U \to X$ in $X_{\et}$, we set
    \[\mu_n(U) = \{ f \in \Gamma(U, \~{O}_U)^\times \mid f^n = 1\}.\]
    Equivalently, $\mu_n$ can be interpreted as the scheme defined by the equation $t^n - 1$, and $\mu_n(U)$ is the group of $n^{th}$ roots of $1$ in $\Gamma(U, \~O_U)$. 
    With these definitions, we obtain an exact sequence of \'etale sheaves
    \[1 \to \mu_n \to \G_m \To{(-)^n} \G_m \to 1,\]
    which is known as the \emph{Kummer sequence}.
\end{ex}

\begin{ex}
    We define presheaves of (not necessarily abelian) groups $\GL_n, \PGL_n: X_{\et}^\op \to \Grp$ by
    \[\GL_n(U) := \GL_n(\Gamma(U,\~O_U)) \quad \text{and} \quad \PGL_n(U) := \Aut\left(\M_n(\~O_U)\right)\]
    for each $U \in X_{\et}$. These are in fact \'etale sheaves on $X$ (this follows from the Skolem-Noether theorem, see Corollary 2.4 in chapter IV of Milne's book \cite{milne_etalebook}). Moreover, we have the following exact sequence of \'etale sheaves on $X$:
    \[1 \to \G_m \to \GL_n \to \PGL_n \to 1\]
\end{ex}


\begin{ex}
Recall that in ordinary sheaf theory, given sheaves $\~F$ and $\~G$ on the same topological space, we can form a sheaf $\sHom(\~F, \~G)$ (``sheaf hom'') of morphisms $\~F \to \~G$. In the same way, given $\~F, \~G \in \Sh(X_{\et})$, we can define a new \'etale sheaf $\sHom_{\et}(\~F, \~G)$ by the assignment
\[U \mapsto \Hom_{\Sh(X_{\et})}(\~F|_U, \~G|_U).\]
\end{ex}

\begin{ex}[The Nisnevich topology]
    An \'etale morphism $f: Y \to X$ is called a \emph{Nisnevich morphism} if for every $x \in X$ there is some $y \in f^{-1}(x)$ for which the induced map of residue fields $\kappa(x) \to \kappa(y)$ is an isomorphism. This is equivalent to the assertion that $f$ is flat, unramified, locally of finite presentation, and for each point $x \in X$ there is $y \in f ^{-1}(x)$ such that $\kappa(x) \to \kappa(y)$ is an isomorphism. A \emph{Nisnevich covering} of $X$ is a family $\{f_i: U_i \to X\}_{i \in I}$ of \'etale morphisms such that for each point $x \in X$ there exists $i \in I$ and $y \in U_i$ such that $f_i(y) = x$ and the induced map $\kappa(x) \to \kappa(y)$ is an isomorphism. If $|I| = n < \infty$, this condition is equivalent to the morphism $\coprod_{i=1}^n U_i \to X$ being Nisnevich. The Nisnevich covers define coverings in a Grothendieck topology, which is unsurprisingly referred to as the \emph{Nisnevich topology}. It is finer than the Zariski topology but coarser than the \'etale topology, and plays a central role in areas like motivic homotopy theory and $\A^1$-enumerative geometry.
\end{ex}

\begin{defn}
    For $n \ge 0$, we define the \emph{$n^{th}$ \'etale cohomology group} of a scheme $X$ by
    \[H^n_{\et}(X, \~F) := R^n\Gamma(X, \~F)\]
    where $\Gamma(X, -): \Sh_{\Ab}(X_{\et}) \to \Ab$ is the global sections functor and $H^0_{\et}(X, \~F) = \Gamma(X,\~F)$.
\end{defn}

With this definition, we obtain for each exact sequence $0 \to \~F \to \~G \to \~H \to 0$ of \'etale sheaves on $X$ a long-exact sequence in \'etale cohomology 
    \[\begin{tikzcd}[sep = scriptsize]
    0 &  H^0_{\et}(X, \~F) &  H^0_{\et}(X, \~G) &  H^0_{\et}(X, \~H) \\
    & { H^1_{\et}(X, \~F)} & { H^1_{\et}(X, \~G)} & { H^1_{\et}(X, \~H) } \\
    & {H^2_{\et}(X, \~F)} & \cdots
    \arrow[from=1-1, to=1-2]
    \arrow[from=1-2, to=1-3]
    \arrow[from=1-3, to=1-4]
    \arrow[from=1-4, to=2-2]
    \arrow[from=2-2, to=2-3]
    \arrow[from=2-3, to=2-4]
    \arrow[from=2-4, to=3-2]
    \arrow[from=3-2, to=3-3]
\end{tikzcd}\]

\subsection{Brauer groups via \'etale cohomology}

So far we have studied Brauer groups in the context of Azumaya algebras, but we've also seen that these groups have a cohomological description. Recall that for a scheme $X$, an $\~O_X$-algebra $A$ is an \emph{Azumaya algebra} over $X$ if it is a coherent $\~O_X$-module and $A_x$ is an Azumaya algebra over the local ring $\~O_{X, x}$ for every point $x$ of $X$. This is equivalent to the assertion that there is an \'etale (resp. flat) covering such that for each $i$ there exists an $d_i \ge 1$ with $A \otimes_{\~O_X} \~O_{U_i} \cong \M_{d_i}(\~O_{U_i})$ (see Proposition 2.1 in Chapter IV of \cite{milne_etalebook}). We will write $\Br(X)$ for the Brauer group of Azumaya algebras over $X$ (as defined in the previous lectures).

\begin{defn}
    The \emph{cohomological Brauer group} of a scheme $X$ is defined by
    \[\Br'(X) := H_{\et}^2(X, \G_m).\]
\end{defn}

\begin{rmk}
    You can also define the cohomological Brauer group in the case of sheaves on the big \'etale site.
\end{rmk}

\begin{thm}
    For a scheme $X$, there is a canonical injection $\Br_{\Az}(X) \mono H^2_{\et}(X, \G_m)$.
\end{thm}
\begin{proof}
    See \cite{milne_etalebook}, Chapter IV, Theorem 2.5.
\end{proof} 

It is natural to ask when this map is surjective, but answering this question proves difficult in general. Here are a few known cases:
\begin{enumerate}
    \item If $X$ is a quasi-projective scheme, then $\Br(X) \cong \Br'(X)$ (cf. \cite{stacks} \href{https://stacks.math.columbia.edu/tag/0A2J}{\S0A2J}).
    \item If $X = \Spec(R)$ with $R$ a Henselian local ring (\cite{stacks}, Definition \href{https://stacks.math.columbia.edu/tag/04GF}{04GF}), we have an isomorphism $\Br(X) \cong \Br'(X)$.
    \item If $X$ is a smooth variety admitting a decomposition $X = X_1 \cup X_2$ with $X_1, X_2$, and $X_1 \cap X_2$ all open affines, then $\Br(X) \cong \Br'(X)$.
    \item A theorem of Gabber (communicated by de Jong in \cite{deJong_Gabber}) says that $\Br(X) \cong \Br'(X)$ for a scheme $X$ admitting an ample invertible sheaf. This implies for example that $\Br(X) \cong \Br'(X)$ when $X$ is affine, since the trivial line bundle on any affine scheme is ample.
\end{enumerate}

\subsection{\'Etale site of $\Spec(k)$ and Galois cohomology}
Fix a field $k$. Our goal in this section is to better understand the relationship between Galois cohomology and \'etale cohomology. This is a nice way of getting acquainted with the small \'etale site, since our attention will be restricted to $\Spec(k)_{\et}$. First, we need to understand what the coverings in this site look like.

\begin{lem}[\cite{stacks}, Lemmas \href{https://stacks.math.columbia.edu/tag/00U3}{00U3} and \href{https://stacks.math.columbia.edu/tag/02GL}{02GL}]\label{lem:etale_in_Speck} \hfill
\begin{enumerate}
\item A morphism of rings $k \to S$ is \'etale if and only if there is an isomorphism of $k$-algebras $S \cong \prod_{i=1}^n k_i$ with $k_i/k$ a finite separable extension for all $i$.
\item A morphism $U \to \Spec(k)$ is \'etale if and only if $U \cong \coprod_{i \in I} \Spec(k_i)$ with $k_i/k$ a finite separable field extension for each $i \in I$.
\end{enumerate}
\end{lem}

% \begin{proof}
%     Applying $\Spec$ to the statement in (1) yields (2), so it suffices to prove (1). If $k'/k$ is a finite separable extension, the primitive element theorem tells us that $k' \cong k(\alpha) \cong k[x]/\langle m_{\alpha} \rangle$, where $m_{\alpha}$ is the minimal polynomial of $\alpha$. As $k'/k$ is separable, we have $\gcd(m_\alpha,m_\alpha') = 1$ (where $m_\alpha'$ is the derivative of $m_\alpha$)
% \end{proof}
This immediately gives the following characterization of \'etale coverings of $\Spec(k)$:
\begin{cor}\label{cor:coverings_in_Speck}
    The coverings in $\Spec(k)_{\et}$ can be written as 
   \[\left(\coprod_{j \in J_i} \Spec(k_{ij}) \to \Spec(k) \right)_{i \in I}\]
   with $k_{ij}/k$ a finite separable extension for all $i,j$.  
\end{cor}
If $g: k' \to k''$  is a morphism between finite separable extension of $k$, we get a morphism $\Spec(g): \Spec(k'') \to \Spec(k')$ in $\Spec(k)_{\et}$, so we have an induced map of sets
\[\~F(g) := \~F\left(\Spec{g}\right) : \~F\left(\Spec{k'}\right) \to \~F\left(\Spec{k''}\right)\]
with $\~F\left(h\right) \circ \~F\left(g\right) = \~F\left(h \circ g\right)$ for $h: k'' \to k'''$. As $\Spec(g)$ is a map in $\Spec(k)_{\et}$, it is a covering map, hence $\~F(g)$ is injective. 

Suppose $k''/k'$ is Galois with corresponding inclusion $\iota: k' \mono k''$. Given an automorphism $\sigma: k'' \to k''$ fixing $k'$, we get a corresponding automorphism $\sigma^*: \Spec(k'') \to \Spec(k'')$ of schemes over $k'$. The group 
\[G := \Gal(k''/k') = \Aut(\Spec{k''}/\Spec{k})^\op\]
acts on $\~F(\Spec(k''))$ via 
\[\sigma \cdot s = \~F(\sigma^*)(s) \quad \text{for all } s \in \~F(\Spec{k''}).\]
Moreover, note that the injection
\[\~F(\iota) : \~F(\Spec{k'}) \mono \~F(\Spec{k''})\] 
is invariant under this action. 

\begin{lem}\label{lem:etSheaves_Speck}
    A presheaf $\~F$ on $\Spec(k)_{\et}$ is an \'etale sheaf (of sets) if and only if the following conditions are satisfied:
    \begin{enumerate}
        \item For any disjoint union $\coprod_{i \in I} U_i$, we have
        \[\~F\left(\coprod_{i \in I} U_i\right) = \prod_{i \in I} \~F(U_i)\]
        \item For all finite separable extensions $k''/k'/k$ with $k''/k'$ Galois, $\~F \left(\Spec{k'}\right) \mono \~F \left(\Spec{k''}\right)$ is injective and $\~F(\Spec{k'}) = \~F(\Spec{k''})^{\Gal(k''/k')}$.
    \end{enumerate}
\end{lem}
\begin{proof}
Now let $U \in (\Spec{k})_{\et}$ and take an \'etale cover $\{U_i \to U\}$. We are interested in the following equalizer diagram:
\[\~F(U) \to \prod_i \~F(U_i) \rightrightarrows \prod_{i,j} \~F(U_i \times_X U_j).\]
Set $V = \coprod_i U_i$. Then we have a surjective \'etale map $V \epi U$, and condition (1) gives $\prod_i \~F(U_i) \cong \~F(V)$. Using the series of isomorphisms
\begin{align*}
\~F(V \times_U V) &\cong \~F\left(\coprod_{i}U_i \times_U \coprod_{j}U_j\right) \\ 
&\cong \~F\left(\coprod_{i,j} (U_i \times_U U_j)\right) \\ 
&\cong \prod_{i,j} \~F(U_i \times_U U_j)
\end{align*}
we can rewrite the sheaf equalizer diagram as follows:
\[\~F(U) \to \~F\left(V\right) \rightrightarrows \~F(V \times_U V).\]
Hence it suffices to check the sheaf axioms for coverings consisting of a single \'etale map $\coprod_i U_i \to U$, which by Corollary \ref{cor:coverings_in_Speck} must have the form $\coprod_{i} \Spec(k_i) \to U$ with $k_{i}/k$ a finite separable extension. Therefore it suffices to check the sheaf condition for a single separable extension $k'/k$, and Galois theory tells us that we can find a field $k''$ with $k''/k'$ Galois. We are therefore reduced to the case of a finite separable extension $k''/k'/k$ with $k''/k'$ Galois. To that end, consider the diagram
\begin{equation}\label{eq:D1}\Spec{k''} \times_{\Spec{k'}} \Spec{k''} \rightrightarrows \Spec{k''} \to \Spec{k'}\end{equation}
where the maps on the left are the projections. Let $G = \Gal(k''/k')$. We have an isomorphism
\begin{align*} 
k'' \otimes_{k'} k'' &\isom \prod_{g \in G} k'' \\
x \otimes y &\mapsto (xg(y))_{g \in G}
\end{align*}
of $k''$ algebras, which allows us rewrite the diagram (\ref{eq:D1}) as follows:
\begin{equation}\label{eq:D2}\coprod_{g \in G} \Spec{k''} \rightrightarrows \Spec{k''} \to \Spec{k'}\end{equation}
This diagram is an equalizer if and only if 
\begin{equation}\prod_{g \in G} \~F(\Spec{k''}) \leftleftarrows \~F(\Spec{k''}) \leftarrow \~F(\Spec{k'})\end{equation}
is an equalizer, where the top-left map is the diagonal morphism $s \mapsto (s, \hdots, s)$ and the bottom-left map is given by $s \mapsto (gs)_{g \in G}$. For this diagram to be an equalizer, the map $\~F(\Spec{k'}) \mono  \~F(\Spec{k''})$ must be injective and $\~F(\Spec{k'}) = \~F(\Spec{k''})^G$.
\end{proof}
\begin{prop}
    $\G_a$ and $\G_m$ are sheaves on $\Spec(k)_{\et}$. 
\end{prop}

\begin{proof}
    It is enough to prove the result for $\G_a$. We need to check that the conditions in Lemma \ref{lem:etSheaves_Speck} are satisfied. Suppose we are given $\coprod_{i \in I}U_i$ with $U_i \in \Spec(k)_{\et}$. Using Lemma \ref{lem:etale_in_Speck}, we can refine this to a disjoint union $\coprod_{i \in I} \Spec(k_i)$ with $k_i/k$ finite separable for all $i \in I$. Then
    \[\G_a\left(\coprod_{i \in I} \Spec{k_i}\right) = \G_a\left(\Spec{\prod_{i \in I} k_i} \right) = \prod_{i \in I} k_i = \prod_{i\in I} \G_a\left(\Spec{k_i}\right),\]
    which gives the first condition in Lemma \ref{lem:etSheaves_Speck}. For the second condition, suppose we have finite separable extensions $k''/k'/k$. The map
    \[\G_a(\Spec{k'}) = k' \to k'' = \G_a(\Spec{k''})\]
    is injective (it is a morphism of fields), and we know from Galois theory that $E^{\Gal(E/F)} = F$, so 
    \[\G_a(\Spec{k''})^{\Gal(k''/k')} = (k'')^{\Gal(k''/k')} = k'= \G_a(\Spec{k'}).\]
    Hence $\G_a$ defines an \'etale sheaf on $\Spec(k)$.
\end{proof}
As for ordinary sheaves, the kernel of a morphism of \'etale sheaves is an \'etale sheaf, and $\mu_n$ is the kernel of the morphism
\begin{align*}(-)^n: \G_m &\to \G_m \\ x &\mapsto x^n\end{align*}
so we get the following corollary:
\begin{cor}
    $\mu_n$ is an \'etale sheaf on $\Spec(k)$.
\end{cor}

\begin{defn}
    A \emph{profinite group} is a topological group $G$ satisfying either of the following equivalent conditions:
    \begin{enumerate}
        \item $G$ is a filtered limit of discrete finite groups: $G \cong \lim_{i \in \mathbf{I}} G_i$ with $\mathbf{I}$ a filtered category and each $G_i$ discrete (put differently, $G$ is a pro-object in the category of finite groups).
        \item $G$ is Hausdorff, compact and totally disconnected (the only connected subsets of $G$ are singletons).
    \end{enumerate}
\end{defn}

\begin{ex}
    Any finite group is a profinite group. If $k$ is a field equipped with a choice of separable closure $k^{\sep}$, then the absolute Galois group $G_k := \Gal(k^{\sep}/k)$ is a profinite group: 
    \[G_k := \varprojlim_{k'/k \text{ finite}} \Gal(k'/k).\]
    Endowing each of the finite groups $\Gal(k'/k)$ with the discrete topology, $G_k$ inherits a topology called the \emph{Krull topology} (it is just important to know that such a topology exists).
\end{ex}

\begin{defn}
If $G$ is a profinite group, a \emph{discrete (left) $G$-set} is a set $M$ equipped with a (left) $G$-action $G \times M \to M$ which is continuous when $M$ is given the discrete topology. This means $G \times M \to M$ is continuous if and only if the stabilizer of every $m \in M$ is an open subgroup of $G$. When $M$ is an abelian group and the $G$-action respects this structure, then we call $M$ a \emph{discrete $G$-module} or \emph{Galois module}. Denote by $\Set_G^\disc$ (resp. $\Mod_G^\disc$) the category of discrete (left) $G$-sets (resp. $G$-modules).
\end{defn}

\begin{rmk}
    The category $\Mod^\disc_{G}$ of discrete $G$-modules is an abelian category with enough injectives, so it is a reasonable setting in which to discuss cohomology (for how to construct injectives in this category, see Remark 1.1.4.2 in \cite{conrad}).
\end{rmk}
For a profinite group $G$ and a short exact sequence of discrete $G$-modules, taking $G$-invariants gives a left exact functor $M \mapsto M^G$: For a short exact sequence 
\[0 \to M_1 \to M_2 \to M_3 \to 0\]
we obtain a left-exact sequence
\[0 \to M_1^G \to M_2^G \to M_3^G.\]
\begin{defn}
For a profinite group $G$ and $M \in \Mod_G^{\disc}$, the \emph{Galois cohomology groups} $H^n(G, M)$ are then defined as the right derived functors of the $G$-invariants functor:
\[H^n(G, M) := R^n(-)^G.\]
\end{defn}
\begin{thm}
    Let $k$ be a field equipped with a choice of separable closure $k^\sep$, and set $G_k = \Gal(k^\sep/k)$. There is an equivalence of categories $\Sh(\Spec(k)_{\et}) \simeq \Set_{G_k}^\disc$, where $\Set_{G_k}^\disc$ denotes the category of left $G_k$-sets.
\end{thm}
\begin{proof}
    Let $\Sigma = \{k' \supseteq k : k'/k \text{ Galois and } k' \subseteq k^\sep\}$. Given $k' \in \Sigma$, we have a set $\~F(\Spec{k'})$ with a left action by $\Gal(k'/k)$ and for $k'' \in \Sigma$ with $k' \subseteq k'' \subseteq k^\sep$, the previous proposition gives an injection $\~F(\Spec{k'}) \mono \~F(\Spec{k'})$ with $\~F(\Spec{k'}) = \~F(\Spec{k''})^{\Gal(k''/k')}$. Moreover, $\~F(\Spec{k'}) \mono \~F(\Spec{k'})$ is compatible with the Galois action via the surjection $\Gal(k''/k) \epi \Gal(k'/k)$. Define
    \[M_{\~F} := \varinjlim_{k' \in \Sigma} \~F(\Spec{k'}).\]
    Then $M_{\~F}$ inherits the structure of a discrete $G_k$-set, and the canonical map $\~F(\Spec{k'}) \to M_{\~F}$ is an isomorphism (of $\Gal(k'/k)$-sets) onto $M_{\~F}^{\Gal(k^\sep/k')}$. The functor $\~F \mapsto M_\~F$ give the desired equivalence. For an explicit proof that this defines an equivalence, see Theorem 1.1.4.3 in \cite{conrad}\footnote{Proving this is \emph{highly} nontrivial; Conrad's proof is over a page long.}.
\end{proof}

\begin{cor}
The right-derived functors of the global sections functor on $\Sh_{\Ab}(\Spec(k)_{\et})$ are exactly the Galois cohomology groups for $k^\sep/k$.
\end{cor}

\begin{prop}[Hilbert 90 for \'etale cohomology]
    We have $H^1_{\et}(\Spec(k), \G_m) = 0$.
\end{prop}
\newpage
\printbibliography
\end{document}
